\section{A constant quadratic enlargement of the signed hosts}
\label{sec:mixed-host-interface}

The fixed cubic coefficient is scalar over a field containing
$K=\mathbb Q(\omega)$.  The literal signed Host in AM-H17 does
not contain this field.  We first make the required enlargement
and its integral cost precise.  Here ``Host'' has exactly the
scope of \texttt{DataHost.lean}; it does not include the later
theta identities or the proposition \texttt{VWin}.

\begin{proposition}[AC-G30: adjoining the cubic coefficient field]
\label{AC-G30}
Write the AM-H17 field as
\[
 C_h=\mathbb Q(\zeta_{2^{h+2}},\zeta_7),\quad
 H=C_hK_qK_d,\quad K_q=\mathbb Q(\sqrt{-q}),\quad
 K_d=\mathbb Q(\sqrt{-d}),
\]
and put $n=3\cdot2^{h+2}$ and $g=[H:\mathbb Q]/2=2n$.
Assume $h\geq294$, the splitting bound is at least $269$,
and the AM-H17 primes are chosen large enough for the ranks
specified below.  Let $F=HK$.  Then
\[
 [F:H]=2,\qquad
 \operatorname{rd}(F)=\sqrt3\operatorname{rd}(H),
 \qquad \mu(F)=\mu_{3\cdot7\cdot2^{h+2}}.             \tag{G30.1}
\]
The AM-H17 construction gives a complete bare Host over $F$,
with the same $q,d$ and
\[
 Q'_S=3Q_S,\quad n'=2n,\quad n'_0=2^h,\quad o'=24.
                                                               \tag{G30.2}
\]
More explicitly, if the original ranks are $(k,m,t)$ and
$t\geq2k+m+4$, take $(k',m',t')=(2k,m,t)$ and require
$q,d>4t(2k+m-1)+1$.  Thus this enlargement preserves the
weight cutoff $t$; these are explicit rank and prime-size
conditions, not a claim that arbitrary old ranks still work.

Take the AM-H17 choice $\mathfrak c_K=\mathcal O_{K_d}$ and
write its different-splitting ideal and lattice as
$\mathfrak b_H$ and $AA_H=\mathfrak b_H^{-1}$.
There is a three-supported ideal $\mathfrak b_3$ of $F$ with
\[
 \mathfrak b_3\overline{\mathfrak b}_3
       =\mathfrak D_{F/H},\quad
 N_{F/\mathbb Q}\mathfrak b_3=3^g,
\]
such that the new source lattice is
\[
 AA_F=(AA_H\mathcal O_F)\mathfrak b_3^{-1},\qquad
 [AA_F:AA_H\mathcal O_F]=3^g.                         \tag{G30.3}
\]
It is unchanged relative to the extended old lattice away
from three.  With the source coefficient
$\xi_0=-1/\sqrt{-d}$ it satisfies the exact principal identity
\[
 (\xi_0)AA_F\overline{AA_F}\mathfrak D_{F/\mathbb Q}=1.
                                                               \tag{G30.4}
\]
Its Euclidean covolume is
\[
 \operatorname{covol}(AA_F)
       =\left(\frac{\sqrt d}{2}\right)^{2g},           \tag{G30.5}
\]
and the determinant of the principal trace metric makes its
metric covolume one.  The logarithm of the integral enlargement
index in (G30.3) is $g\log3$; that logarithmic cost must not be
substituted for an unlogged theta norm.

Nevertheless neither $H$ nor $F$ is a literal
\texttt{IsChiField} field for any one Dirichlet character
with these same $h,q,d$.  This last assertion concerns the
single-character definition; it does not assert that several
instances of a separately supplied theorem \texttt{VWin}
could never imply a bound for these fields.
\end{proposition}
\begin{proof}
\proofstep{1. The field and frame.}
The four factors used in AM-H17 have ramification support
in $\{2\}$, $\{7\}$, $\{q\}$ and $\{d\}$.
The primes $q,d$ exceed $269$.  Thus $H$ is unramified at
three and $H\cap K=\mathbb Q$, since $K/\mathbb Q$
ramifies at three.  The discriminants are coprime, so
\[
 \mathcal O_F=\mathcal O_H\otimes_{\mathbb Z}\mathcal O_K,
 \qquad
 |\operatorname{disc}F|
       =|\operatorname{disc}H|^2\,3^{[H:\mathbb Q]}.
\]
This is the integral-basis/discriminant formula used in
AM-H17, with the additional quadratic factor of discriminant
$-3$; see the primary locator there, Milne,
\emph{Algebraic Number Theory}, Lemma~6.5 and Remark~6.6(c),
printed pp.~100--101.  Taking the $[F:\mathbb Q]$-th root
gives (G30.1).  The roots-of-unity argument in AM-H17 applies
at every old prime, while the exact inertia index at three
is two and excludes a primitive ninth root.  It gives the
displayed root group.

Use the frame group
\[
 A'=\operatorname{Gal}(C_hK/\mathbb Q)
        \simeq C_2\times C_{2^h}\times C_6\times C_2.
\]
Project to the two-power cyclotomic factor modulo conjugation,
and use the element $5$ in that factor as the distinguished
generator, acting trivially on the other factors.  Its image
has order $2^h$, with the same injectivity on the distinguished
cyclic subgroup as in AM-H17.  Thus $|A'|=2n$ and
$o'=|A'|/2^h=24$.  The integer inequality
$12^{64}<2^{230}$ used in AM-H17 gives
$24^{64}<2^{294}\leq2^h$.

All rank clauses in \texttt{Params.Ranks} hold for
$(2k,m,t)$ under the stated extra slack.  In particular
$2n'=4n\leq2k$, the parity conditions are unchanged,
and $k'+m-1+5\leq t$ is exactly the added condition.
The remaining inequalities either stay the same or become
weaker when $k$ increases.  The strengthened lower bound
on $q,d$ is exactly the new \texttt{qd\_big} clause.

\proofstep{2. All bare Host field clauses survive.}
Use the rebasing map of AM-H17 with the conjugation of
$C_hK$ in place of that of $C_h$.  It again sends full
complex conjugation to $((-1,-1),1)$ and pure $d$-inertia
to $((1,-1),1)$.  Restriction from the cyclotomic field of
conductor $3Q_Sqd$ supplies the embedding, Artin map,
quadratic-coordinate identities, and Frobenius, inertia
and decomposition descriptions in \texttt{HostRaw} and
\texttt{HostDer}.  The new conductor is exact at three
because its quadratic factor acts nontrivially; exactness
at the old primes is preserved by fixing the added factor.
The inequalities $Q'_S\leq(n')^2$ and the required
root-discriminant bounds follow from
$Q'_S=3Q_S$, $n'=2n$, and (G30.1).

There is only one new local check.  Three splits in $K_d$
by the stipulated small-prime splitting condition.  Hence
its decomposition group fixes $K_d$ and does not contain
complex conjugation.  The old ramified-prime checks remain
valid: if conjugation belonged to a new decomposition
group, its restriction would belong to the old one.
Thus the CM extension $F/F^+$ is unramified at every finite
place, and the different at the new prime can be separated
into conjugate halves.  The old proofs of the ordinary-type,
quadratic splitting and sign-character clauses use the same
$K_d$ and the same two quadratic coordinates, so they apply
without alteration to the added factor.

The old four auxiliary primes in the AM-H17 construction
were chosen congruent to one modulo $15Q_Sqd$.
They are therefore already congruent to one modulo three
and split in the added quadratic factor.  They retain all
bare Host level and coprimality requirements.  One can keep
$\mathfrak c_K=\mathcal O_{K_d}$, and the choices of empty
$TT,T_2$ allowed for the bare structure.  As in AM-H17, this
last observation does not claim the later quantitative
lower bounds on these sets.

The exact declarations checked here are
\texttt{HostRaw}, \texttt{HostDer}, \texttt{HostAx} in
\path{Formalism/SharpRestriction/RamificationFrobenius/TotientGaloisClasses/ConjEMemHostFIff.lean},
and \texttt{Host} in \texttt{DataHost.lean}.
The separately defined property
\texttt{HostRaw.CoprimeFifteen} is not a field of this Host;
the new conductor is divisible by three and does not satisfy
that separate property.

\proofstep{3. The exact principal lattice at three.}
Since the discriminants are coprime,
\[
 \mathfrak D_{F/H}=(\sqrt{-3})\mathcal O_F.
\]
Choose one of the two primes of $K_d$ above three, and let
$\mathfrak b_3$ be the part of this relative different above
that chosen prime.  Conjugation exchanges the two supports,
and every automorphism fixing $K_d$ preserves each support.
Consequently
\[
 (\mathfrak b_3,\overline{\mathfrak b}_3)=1,
 \qquad
 \mathfrak b_F=\mathfrak b_H\mathcal O_F\,\mathfrak b_3
\]
is the required stable half of $\mathfrak D_{F/K_d}$.
Its inverse is (G30.3).  The norm of the relative different
is $3^{2g}$; its two conjugate halves have equal norm, so
$N\mathfrak b_3=3^g$.  For invertible fractional ideals,
$[I\mathfrak b_3^{-1}:I]=N\mathfrak b_3$, which proves
the index.  Transitivity of differents and the old principal
identity give (G30.4) directly.

The extended old lattice has Euclidean covolume
\[
 2^{-2g}\sqrt{|\operatorname{disc}F|}\,N(AA_H)^2
       =3^g\left(\frac{\sqrt d}{2}\right)^{2g}.
\]
Division by its enlargement index gives (G30.5).
The old positive trace metric has diagonal coefficient
$2/\sqrt d$ in each of the $2g$ complex coordinates, so
its volume density cancels this covolume.  For the signed
metric obtained by pulling back the real unit $\kappa$,
the product of $|\tau(\kappa)|$ remains one, and the
same determinant calculation applies.  These statements
compare the actual trace lattices and determinant metrics;
they do not identify normalized algebraic theta sections.

Finally the roots of unity in $F$ are those in $C_hK$.
Their norm to $K_d$ equals their rational norm from that
CM cyclotomic factor, squared, and is one.  This verifies
the remaining \texttt{norm\_mu} clause of the bare Host.

\proofstep{4. The single-character obstruction.}
Put
\[
 B_h=\mathbb Q(\zeta_{2^{h+2}})^+K_qK_d.
\]
The literal definition \texttt{IsChiField} in
\path{Development/AutomorphicArithmetic/CharacterBoxes/ExistsPageN.lean}
requires its field to be
\[
                           B_hL_\psi,                 \tag{G30.6}
\]
where $L_\psi$ is the cyclic field cut out by the one finite
Dirichlet character $\psi$.  Indeed its displayed stabilizer
is the intersection of the kernel of $\psi$ with the three
stabilizers defining the real two-power field and the two
quadratic fields.  Galois correspondence identifies this
intersection with the compositum in (G30.6), regardless
of which sufficiently divisible ambient conductor is used.
The finite character image is cyclic, so
$\operatorname{Gal}(B_hL_\psi/B_h)$ is cyclic.

By disjoint ramification of the factors,
\[
 \operatorname{Gal}(H/B_h)=C_2\times C_6,
 \qquad
 \operatorname{Gal}(F/B_h)=C_2\times C_6\times C_2.
                                                               \tag{G30.7}
\]
Neither group is cyclic.  Thus neither field can satisfy
the same single-character definition.  This proves the
precise obstruction claimed, while leaving open a comparison
using several cyclic subfields and several quantitative
instances of \texttt{VWin}.
\end{proof}

\begin{remark}[The scalar coefficient and the one-Speh-block gate]
Over $H$, the splitting field of $X^3-269$ still has group
$S_3$: its only possible nontrivial abelian intersection
with the abelian field $H$ is $K$, and $K\not\subset H$.
The rational rank-two Artin representation therefore remains
irreducible over $H$, even after extending its coefficient
field.  Over $F=HK$ it splits into the two cubic characters.
The CM-linear Hom construction of AC-G23 is scalar over
this latter base.

On $F$, scalar twists of separate packets preserve their
inertia types at the two old auxiliary primes, since the
cubic character is unramified there.  This avoids the
repeated-inertia-polynomial defect of the full regular
tensor in AC-G29.  It does not repair the source's
arithmetic interface.  The exact declarations
\texttt{NData.chi\_def}, \texttt{muPart},
\texttt{HasParamX}, \texttt{ZN} and \texttt{IsRhoBar}
in \texttt{LevelLetterCocycles/DataN.lean} use the one
distinguished Speh block
$\langle1,\mu_C,t\rangle$ and its characters
$\eta_c=\eta_1\varepsilon_\ell^{-(c-1)}$.
All pairwise ratios inside that block are pure Tate.

For example, replacing these characters by
$\eta_c=\eta_0\chi^{i_c}\varepsilon_\ell^{-(c-1)}$
can satisfy that same \texttt{chi\_def} only if all
$i_c$ agree modulo three.  Evaluate the ratio at the
three-primary tame inertia at $269$: $\chi$ is surjective,
whereas $\varepsilon_\ell$ is trivial there.  For
$\ell\ne269$ this is the usual unramified cyclotomic
action; for $\ell=269$ it follows because
$\mathbb Z_{269}^\times$ has no nontrivial pro-three
subgroup.  Thus $\omega^{i_c-i_1}=1$ for every $c$.
Separate scalar packets consequently require a new
comparison for extensions between distinct distinguished
blocks, together with its integral theta and Hecke ties.
The bare Host enlargement above supplies none of those
additional identities.
\end{remark}

\begin{proposition}[AC-G30-Fourier: fixed finite-character Fourier defects]
\label{AC-G30-Fourier}
For the extension $F/B_h$ in (G30.7), let
$G=\operatorname{Gal}(F/B_h)\simeq C_2^3\times C_3$.
Its characters are the restrictions of the $24$ Dirichlet
characters of modulus $84$.  Let
\[
 \mathcal F:\mathcal O[G]\longrightarrow
       \bigoplus_{\eta\in\widehat G}\mathcal O(\eta)
\]
be the integral Fourier map.
At three, take $\mathcal O=\mathbb Z_3[\omega]$.
The cokernel is killed by three and has
$\mathcal O$-length $24$.
At two, take the unramified quadratic ring
$\mathcal O=\mathbb Z_2[\omega]$.
The elementary divisors are
\[
 1\ (3\text{ times}),\qquad
 2\ (9\text{ times}),\qquad
 4\ (9\text{ times}),\qquad
 8\ (3\text{ times}).                                \tag{G30.8}
\]
Thus this cokernel is killed by eight and has
$\mathcal O$-length $36$, or underlying
$\mathbb Z_2$-length $72$.
At all other coefficient primes the Fourier map is an
integral isomorphism after an unramified coefficient
extension containing the sixth roots of unity.

These are fixed finite defects, independent of $h,q,d$
and of a common Tate twist.  Their residual characters
need not be trivial.  The common fixed extension
$\mathbb Q(\zeta_{168})$ trivializes their finite actions
and the required finite Tate twists after base change
from $B_h$.
\end{proposition}
\begin{proof}
The full two-power cyclotomic field is obtained from its
maximal real subfield by adjoining $i$.  Consequently
\[
 F=B_h\mathbb Q(i,\zeta_7,\omega)
                  =B_h\mathbb Q(\zeta_{84}).
\]
Its relative degree is $24$ by (G30.7), equal to
$\varphi(84)$, so restriction identifies the two Galois
groups.  This proves the character assertion.

Fourier evaluation on a product of groups is the tensor
product of their evaluation matrices.  Over any ring in
which two is not a unit, the $C_2$ matrix
$\left(\begin{smallmatrix}1&1\\1&-1\end{smallmatrix}\right)$
has Smith factors $1,2$, by integral row and column
operations.  The threefold tensor product has factors
$2^j$ with multiplicities $\binom3j$.  At two the $C_3$
Fourier matrix is invertible, since three is a unit, and
adds three copies of each factor.  This gives (G30.8).
The residue field of $\mathbb Z_2[\omega]$ has degree two,
which gives the stated underlying length.

At three all three $C_2$ matrices are invertible.
AM-H06 computes the $C_3$ matrix factors
$1,\lambda,\lambda^2$, of length three and killed by
$\lambda^2$, a unit times three.  Eight copies give the
length $24$.  If the coefficient prime divides neither
two nor three, every group order is a unit, and the usual
inverse Fourier matrix is integral.

Twisting both lattices by a common character changes
none of their underlying Smith factors.  Their untwisted
Galois actions factor through $\mathbb Q(\zeta_{84})$.
The three-primary quotient is killed by three, and the
two-primary quotient by eight; adjoining $\mu_3$ and
$\mu_8$ therefore kills every residual Tate action as
well.  The single field $\mathbb Q(\zeta_{168})$ contains
these fields.  This last argument retains the nontrivial
prime-to-coefficient characters, rather than declaring
every composition factor trivial on the original base.
\end{proof}

\section{What descent actually preserves in the source windows}
\label{sec:finite-window-descent}

\begin{proposition}[AC-G31: exact two-power descent and its window cost]
\label{AC-G31}
For a number field $L$ and $e\geq1$, let
$S_\ell(L,e)$ be the strict group with coefficient
$\mathbb Q_\ell/\mathbb Z_\ell(-e)$: it is unramified
outside $\ell$ and zero at each completion above $\ell$.
Suppose $L'/L$ is cyclic of degree $2^s$ and all the
groups over $L'$ in the range $1\leq e<t$ are finite.
Restriction maps strict groups to strict groups, is
injective for odd $\ell$, and satisfies
\[
 \operatorname{length}_{\mathbb Z_2}
       \ker\bigl(S_2(L,e)\longrightarrow S_2(L',e)\bigr)
                  \leq s.                            \tag{G31.1}
\]
In particular a bound \texttt{VVle}$(L',t,A)$ implies
the mathematical bound
\[
 \texttt{VVle}\bigl(L,t,A+s(t-1)\log2\bigr).         \tag{G31.2}
\]
Here the groups are the actual strict Tate groups identified
by the source's \texttt{LData}; no Artin coefficient has
been inserted into that definition.

Fix a finite set of Dirichlet characters $\psi_i$, a single
$0<\epsilon<1/52$, and the monotone page function
\[
 a(x)=\left\lceil(2+8\epsilon)\log_2 x\right\rceil
                         \quad(x\geq1).
\]
Assume the corresponding source propositions
\texttt{VWinAt}$(\psi_i,401)$ have been supplied, with
window constants $J_i$.  Put $J_* =\max(2,\max_i J_i)$.
For all sufficiently large $Z$, and any one pair $q,d$
satisfying \texttt{WinPair}$(\epsilon,Z,q,d)$, let
$L_i(Z)$ be its cyclic \texttt{IsChiField} at page $a(Z)$.
Then each such field satisfies, with
$t_Z=\lfloor Z^{401}\rfloor$,
\[
 \texttt{VVle}\left(L_i(Z),t_Z,
   \exp\left(\tfrac12 Z^{\epsilon J_*/2}\right)
      +(t_Z-1)\left((2+8\epsilon)(J_*/2-1)\log Z
                                      +\log2\right)\right).
                                                               \tag{G31.3}
\]
This is a common lower-page comparison.  It does not give
the required source budget
$\exp(\tfrac12 Z^\epsilon)$ without further control of
the window constants or a simultaneous-window theorem.
\end{proposition}
\begin{proof}
Restriction preserves the local zero conditions and the
unramified conditions: each new inertia group maps into
the old inertia group, even at places ramified in $L'/L$.
Inflation--restriction embeds the kernel of global
restriction into
\[
 H^1\left(\operatorname{Gal}(L'/L),
  (\mathbb Q_\ell/\mathbb Z_\ell(-e))^{G_{L'}}\right).
                                                               \tag{G31.4}
\]
For odd $\ell$, the group order is a unit on the
$\ell$-primary coefficient module, so averaging makes
this group zero.  For $\ell=2$, cyclic cohomology expresses
(G31.4) as a subquotient of the indicated invariant group
and annihilates it by $2^s$.  The invariant group is a
subgroup of the rank-one group $\mathbb Q_2/\mathbb Z_2$.
Every bounded-exponent subquotient of such a group is
cyclic; if killed by $2^s$ it has order at most $2^s$.
This proves (G31.1).  Finiteness descends, and summing the
length bounds for $1\leq e<t$, with the prime weight
$\log\ell$, gives (G31.2) for every finite prime set.
The universal choice of \texttt{LData} in that definition
causes no extra condition: its set of places is fixed,
its character pulls back to the specified cyclotomic
character, and changing a chosen place in the algebraic
closure conjugates the local restriction map.  These
changes preserve the strict kernels and their lengths.

For each $i$ apply its \texttt{VWinAt} at
$X_i=Z^{J_i}$.  Its \texttt{WinPair} input is exactly
the common pair at $X_i^{1/J_i}=Z$.
It supplies a number $Y_i$ with
\[
 Z\leq Y_i\leq Z^{J_i/2}\leq Z^{J_*/2}
\]
and the \texttt{VAt} bound at $x=Y_i$.
If a stated $J_i<2$ occurs, these inequalities exclude
such a pair for $Z>1$; there is no need to assume its
existence in that vacuous case.
The supplied bound is
$\exp(\tfrac12Y_i^\epsilon)$ for weights up to
$\lfloor Y_i^{401}\rfloor$, which includes $t_Z$.

By the compositum description (G30.6), increasing the page
from $a(Z)$ to $a(Y_i)$ extends only the real two-power
cyclotomic factor.  Thus
$L_i(Y_i)/L_i(Z)$ is cyclic of degree $2^{s_i}$, with
\[
 s_i\leq a(Y_i)-a(Z)
 \leq(2+8\epsilon)(J_*/2-1)\log_2Z+1.
\]
Apply (G31.2) and use
$Y_i^\epsilon\leq Z^{\epsilon J_*/2}$ to obtain
(G31.3).  No intersection of the different supplied
windows has been assumed.
\end{proof}

\begin{remark}[The remaining quantitative condition]
The exact source definitions are \texttt{VVle},
\texttt{IsChiField}, \texttt{WinPair}, \texttt{PageN},
\texttt{VAt} and \texttt{VWinAt} in
\texttt{CharacterBoxes/ExistsPageN.lean}.
They use $t=\lfloor x^{401}\rfloor$ and budget
$\sqrt{\exp(x^\epsilon)}$.
In (G31.3), the additional restriction-kernel term is
polynomial in $Z$ and $\log Z$ and would be negligible
relative to that budget.  The loss in the exponent of the
first term is the actual unresolved issue.

Replacing $\epsilon$ by a smaller value does not by itself
resolve it: the supplied $J_i$ depends on that value,
and the definition gives no upper control on
$\epsilon J_i(\epsilon)$.  There is also a second
constraint on naive reparameterization.  Changing
$Z$ to $Z^{\epsilon/\epsilon'}$ preserves the scale
$Z^{\epsilon/5}$ that bounds $\log q$, but it multiplies
the central scale for $\log d$ by
$(\epsilon/\epsilon')^2$, because \texttt{WinPair}
also contains $(\log Z)^2$.
If that multiplier exceeds two, the two prescribed
factor-two intervals for $\log d$ are disjoint.
Thus neither independent windows nor a change of epsilon
has been promoted to a simultaneous supply theorem.
Finite character decompositions and their integral errors
could still be used after such a theorem is proved.
\end{remark}
